\subsection{Full block diagram}
\label{sec:fullblock}

Figure~\ref{fig:full} shows the CPU and the system in one diagram. The diagram shows the data paths. It does not show these items:

\begin{itemize}
  \item The pipeline registers between the stages. Section~\ref{sec:pipeline} shows them.
  \item The control signals from the instruction decoder to the multiplexers and the units.
  \item The AMO unit in MEM and the store lane logic. Section~\ref{sec:alu} shows them.
\end{itemize}

The three horizontal buses in the CPU are the operand buses: rs1, rs2 and the immediate. A dot shows a connection to a bus. A line that crosses a bus without a dot is not connected to it.

\begin{landscape}
\begin{figure}[p]
  \centering
  \begin{tikzpicture}[x=1cm, y=1cm, font=\scriptsize,
      dot/.style={circle, fill, inner sep=1.3pt},
      rail/.style={double, double distance=1.8pt, line width=0.5pt}]

    % ================================================================ CPU
    \draw[dashed, line width=0.6pt] (0,0.35) rectangle (9.9,14.75);
    \node[align=center, font=\small] at (8.6,1.3) {CPU core\\\src{rtl/cpu.vhd}};

    % ---- fetch
    \node[mux, minimum height=1.4cm, minimum width=0.5cm] (pcmux) at (1.6,13.5) {};
    \node[below, font=\tiny] at (2.55,13.45) {PC MUX};
    \node[pblk, minimum width=1.8cm, minimum height=0.6cm] (pc) at (4.2,13.5) {Program counter};
    \draw[pbus] (1.85,13.5) -- (pc.west);
    \node[pblk, minimum width=0.7cm, minimum height=0.36cm] (p4) at (3.2,14.35) {$+4$};
    \draw[rail] (5.6,13.5) -- (5.6,14.35) -- (p4.east);
    \draw[pbus] (p4.west) -- (0.9,14.35) -- (0.9,13.95) -- (1.35,13.95);

    \node[pblk, minimum width=1.8cm, minimum height=0.7cm] (pd) at (8.6,12.3) {Predecode};
    \foreach \i in {0,...,3}
      \node[pblk, minimum width=1.4cm, minimum height=0.28cm, font=\tiny] (st\i) at (6.3,{12.07-0.28*\i}) {STACK \i};
    \draw[pbi] (7.0,12.08) -- (7.7,12.08);
    \draw[pbus] (7.7,12.45) -- (0.9,12.45) -- (0.9,13.05) -- (1.35,13.05);
    \node[above] at (3.6,12.45) {predicted target};

    \node[pblk, minimum width=1.8cm, minimum height=0.8cm] (iq) at (8.6,10.9) {Instruction\\queue (2)};
    \draw[pbus] (pd.south) -- (iq.north);
    \node[pblk, minimum width=2.6cm, minimum height=1.0cm] (dec) at (8.4,9.3) {Instruction\\decode \& control};
    \draw[pbus] (8.6,10.5) -- (8.6,9.8);

    % ---- register file and forwarding
    \node[pblk, minimum width=2.4cm, minimum height=1.3cm] (rf) at (2.3,9.3) {Register file\\32 $\times$ 32 bits\\(M4K RAM)};
    \draw[pbus] (dec.west) -- (rf.east);
    \node[above] at (5.3,9.3) {rs1, rs2, rd};
    \bw[below]{5.3,9.3}{3 $\times$ 5}
    \draw[line width=0.6pt, fill=white] (1.0,7.95) -- (3.6,7.95) -- (3.3,7.45) -- (1.3,7.45) -- cycle;
    \node at (2.3,7.7) {Forward MUX};
    \draw[pbus] (1.8,8.65) -- (1.8,7.95);
    \draw[pbus] (2.8,8.65) -- (2.8,7.95);

    % ---- operand buses (drawn first, so later lines cross over them)
    \draw[rail] (1.15,6.7) -- (7.8,6.7);
    \draw[rail] (2.3,6.4) -- (9.0,6.4);
    \draw[rail] (3.0,6.1) -- (8.4,6.1);
    \node[right] at (7.85,6.78) {rs1};
    \node[right] at (5.1,5.82) {imm};
    \draw[rail] (1.8,7.45) -- (1.8,6.7);   \node[dot] at (1.8,6.7) {};
    \draw[rail] (2.8,7.45) -- (2.8,6.4);   \node[dot] at (2.8,6.4) {};
    \draw[rail] (8.4,8.8) -- (8.4,6.1);    \node[dot] at (8.4,6.1) {};

    % ---- execute units
    \draw[line width=0.6pt, fill=white] (1.9,5.62) -- (3.4,5.62) -- (3.1,5.18) -- (2.2,5.18) -- cycle;
    \node[font=\tiny] at (2.65,5.4) {MUX};
    \draw[pbus] (2.3,6.4) -- (2.3,5.62);   \node[dot] at (2.3,6.4) {};
    \draw[pbus] (3.0,6.1) -- (3.0,5.62);   \node[dot] at (3.0,6.1) {};
    \draw[line width=0.6pt, fill=white] (0.7,4.6) -- (1.6,4.6) -- (1.9,4.35) -- (2.2,4.6) -- (3.1,4.6) -- (2.6,3.8) -- (1.2,3.8) -- cycle;
    \node at (1.9,4.05) {ALU};
    \draw[pbus] (1.15,6.7) -- (1.15,4.6);  \node[dot] at (1.15,6.7) {};
    \draw[pbus] (2.65,5.18) -- (2.65,4.6);

    \node[pblk, minimum width=1.5cm, minimum height=0.9cm] (md) at (4.4,4.4) {Multiply\\divide};
    \draw[pbus] (4.1,6.7) -- (4.1,4.85);   \node[dot] at (4.1,6.7) {};
    \draw[pbus] (4.7,6.4) -- (4.7,4.85);   \node[dot] at (4.7,6.4) {};
    \node[pblk, minimum width=1.4cm, minimum height=0.9cm] (br) at (6.2,4.4) {Branch\\compare};
    \draw[pbus] (5.9,6.7) -- (5.9,4.85);   \node[dot] at (5.9,6.7) {};
    \draw[pbus] (6.5,6.4) -- (6.5,4.85);   \node[dot] at (6.5,6.4) {};
    \node[pblk, minimum width=1.5cm, minimum height=0.9cm] (ad) at (8.1,4.4) {Address\\adder};
    \draw[pbus] (7.8,6.7) -- (7.8,4.85);   \node[dot] at (7.8,6.7) {};
    \draw[pbus] (8.4,6.1) -- (8.4,4.85);

    % ---- CSRs, traps, result
    \node[pblk, minimum width=2.2cm, minimum height=1.0cm] (csr) at (6.3,2.6) {CSRs, traps,\\next PC};
    \draw[pwire] (6.2,3.95) -- (6.2,3.1);
    \node[pblk, minimum width=1.5cm, minimum height=0.5cm] (la) at (8.3,3.5) {Load align};
    \draw[line width=0.6pt, fill=white] (1.4,2.85) -- (4.8,2.85) -- (4.5,2.35) -- (1.7,2.35) -- cycle;
    \node at (3.0,2.6) {Result MUX};
    \draw[pbus] (1.9,3.8) -- (1.9,2.85);
    \draw[pbus] (3.9,3.95) -- (3.9,2.85);
    \draw[pbus] (la.west) -- (4.45,3.5) -- (4.45,2.85);
    \draw[pbus] (csr.west) -- (4.65,2.6);
    % result bus: back to the register file and the forward MUX
    \draw[pbus] (3.1,2.35) -- (3.1,1.25) -- (0.65,1.25) -- (0.65,9.3) -- (rf.west);
    \draw[pbus] (0.65,7.7) -- (1.15,7.7);  \node[dot] at (0.65,7.7) {};
    \node[above] at (1.9,1.25) {result bus};
    \bw[below]{2.4,1.25}{32}
    % redirect: to the PC MUX
    \draw[pbus] (6.3,2.1) -- (6.3,0.75) -- (0.3,0.75) -- (0.3,13.5) -- (1.35,13.5);
    \node[above] at (4.9,0.75) {redirect target};
    \bw[below]{4.2,0.75}{32}

    % ================================================================ caches and bus
    \node[pblk, minimum width=1.8cm, minimum height=2.0cm] (ic) at (12.2,12.9) {I-cache\\8 KB\\512 lines};
    \node[pblk, minimum width=1.8cm, minimum height=3.4cm] (dc) at (12.2,4.9) {D-cache\\8 KB\\512 lines\\[4pt]write buffer\\4 entries};
    \draw[pbus] (pc.east) -- (11.3,13.5);
    \node[above] at (7.6,13.5) {fetch address};
    \bw[above]{10.65,13.5}{32}
    \draw[pbus] (11.3,12.3) -- (pd.east);
    \node[below, fill=white, inner sep=1.5pt] at (10.62,12.12) {instruction};
    \bw[above]{10.65,12.3}{32}
    \draw[pbus] (9.0,6.4) -- (11.3,6.4);
    \node[below] at (10.6,6.3) {store data};
    \node[above] at (9.35,6.42) {rs2};
    \bw[above]{10.65,6.4}{32}
    \draw[pbus] (ad.east) -- (11.3,4.4);
    \node[below] at (10.6,4.3) {address};
    \bw[above]{10.65,4.4}{32}
    \draw[pbus] (11.3,3.5) -- (la.east);
    \node[below] at (10.6,3.4) {load data};
    \bw[above]{10.65,3.5}{32}

    \node[pblk, minimum width=1.0cm, minimum height=10.1cm] (bus) at (14.3,8.95) {\rotatebox{90}{membus (one transaction at a time)}};
    \draw[pbi] (ic.east) -- (13.8,12.9);
    \draw[pbi] (dc.east) -- (13.8,4.9);

    % ================================================================ slaves
    \node[pblk, minimum width=2.2cm, minimum height=0.8cm] (rom)   at (16.7,13.4) {Boot ROM\\4 KB};
    \node[pblk, minimum width=2.2cm, minimum height=1.1cm] (sram)  at (16.7,11.5) {SRAM\\controller};
    \node[pblk, minimum width=2.2cm, minimum height=1.1cm] (sdram) at (16.7,8.6)  {SDRAM\\controller};
    \node[pblk, minimum width=2.2cm, minimum height=1.1cm] (flash) at (16.7,6.6)  {Flash\\controller};
    \node[pblk, minimum width=2.2cm, minimum height=4.8cm] (io)    at (16.7,3.2)  {I/O registers\\[3pt]timers\\CLINT, PLIC\\16550\\syscon\\keyboard FIFO};
    \foreach \y/\a in {13.4/0x0, 11.5/0x1, 8.6/0x2, 6.6/0x3, 4.6/0xF} {
      \draw[pbi] (14.8,\y) -- (15.6,\y);
      \node[above, font=\tiny\ttfamily, inner sep=2.5pt] at (15.2,\y) {\a};
    }
    \node[below, font=\tiny\ttfamily, inner sep=2.5pt] at (15.2,8.6) {0x8};

    % ================================================================ device blocks and pins
    \node[pblk, minimum width=1.6cm, minimum height=0.9cm] (vid)  at (19.1,10.0) {Video};
    \node[pblk, minimum width=1.6cm, minimum height=0.7cm] (uart) at (19.1,5.0)  {UART};
    \node[pblk, minimum width=1.6cm, minimum height=0.7cm] (ps2)  at (19.1,3.9)  {PS/2 receiver};
    \node[pblk, minimum width=1.6cm, minimum height=0.8cm] (jtag) at (19.1,2.8)  {JTAG loader\\and console};
    \draw[pbi] (17.4,10.95) -- (17.4,10.0) -- (vid.west);
    \node[below, font=\tiny] at (17.75,9.95) {video port};
    \draw[pbi] (17.8,5.0) -- (uart.west);
    \draw[pbus] (ps2.west) -- (17.8,3.9);
    \draw[pbi] (17.8,2.8) -- (jtag.west);

    % pins: \y, label, kind (o = output, i = input, b = both), source x
    \foreach \y/\t/\k/\x in {
        11.9/{SRAM\_ADDR[17:0]}/o/17.8, 11.5/{SRAM\_DQ[15:0]}/b/17.8, 11.1/{SRAM control}/o/17.8,
        10.2/{VGA\_R, G, B [3:0]}/o/19.9, 9.8/{VGA\_HS, VGA\_VS}/o/19.9,
        9.0/{DRAM\_ADDR, BA}/o/17.8, 8.6/{DRAM\_DQ[15:0]}/b/17.8, 8.2/{DRAM control, CLK}/o/17.8,
        7.0/{FL\_ADDR[21:0]}/o/17.8, 6.6/{FL\_DQ[7:0]}/b/17.8, 6.2/{FL control}/o/17.8,
        5.2/{UART\_TXD}/o/19.9, 4.8/{UART\_RXD}/i/19.9,
        4.1/{PS2\_CLK}/i/19.9, 3.7/{PS2\_DAT}/i/19.9,
        2.8/{USB-Blaster}/b/19.9,
        1.9/{LEDR, LEDG}/o/17.8, 1.5/{HEX3 to HEX0}/o/17.8, 1.1/{SW[9:0], KEY[3:1]}/i/17.8} {
      \node[pin] (p) at (20.6,\y) {};
      \node[right, font=\tiny\ttfamily, inner sep=2pt] at (20.72,\y) {\t};
      \if\k o\draw[pwire] (\x,\y) -- (p.west);\fi
      \if\k i\draw[pwire] (p.west) -- (\x,\y);\fi
      \if\k b\draw[line width=0.5pt, {Latex[length=1.6mm]}-{Latex[length=1.6mm]}] (\x,\y) -- (p.west);\fi
    }

    % ================================================================ reset, clock, interrupts
    \node[pblk, minimum width=2.4cm, minimum height=0.9cm] (rst) at (12.6,1.3) {Reset and clock\\50 MHz};
    \node[pin] (pck) at (11.7,0.3) {};
    \node[right, font=\tiny\ttfamily, inner sep=2pt] at (11.82,0.3) {CLOCK\_50};
    \node[pin] (pk0) at (13.5,0.3) {};
    \node[right, font=\tiny\ttfamily, inner sep=2pt] at (13.62,0.3) {KEY[0]};
    \draw[pwire] (pck.north) -- (11.7,0.85);
    \draw[pwire] (pk0.north) -- (13.5,0.85);
    \draw[pwire] (rst.west) -- (9.9,1.3);
    \node[above, font=\tiny] at (10.65,1.3) {reset, clk};
    \draw[pwire] (15.6,2.4) -- (csr.east |- 0,2.4);
    \node[above, font=\tiny] at (12.8,2.4) {interrupts: MTIP, MSIP, MEIP};
    \draw[pwire] (15.6,1.3) -- (rst.east);
    \node[above, font=\tiny] at (14.7,1.3) {reboot};

    % ================================================================ legend
    \draw[pbus] (18.5,14.4) -- (19.3,14.4);
    \node[right] at (19.4,14.4) {bus, width in bits};
    \draw[pwire] (18.5,13.95) -- (19.3,13.95);
    \node[right] at (19.4,13.95) {single signals};
    \node[pin] at (18.9,13.5) {};
    \node[right] at (19.4,13.5) {FPGA pin};
    \node[dot] at (18.9,13.05) {};
    \node[right] at (19.4,13.05) {connection to a bus};
  \end{tikzpicture}
  \caption{The full block diagram}
  \label{fig:full}
\end{figure}
\end{landscape}
